\part{Perspectiva Teórica}

\chapter*{Marco conceptual}

En el ámbito de la gestión de proyectos y el desarrollo de software, es esencial contar con marcos
de referencia que ofrezcan una estructura robusta y metodologías efectivas para enfrentar los
retos actuales. Este proyecto se fundamenta en varios marcos de referencia reconocidos a nivel
internacional, cada uno con su enfoque particular y conjunto de principios que orientan las
prácticas organizacionales y técnicas. Entre estos se incluyen el PMBOK, que establece buenas
prácticas en la gestión de proyectos; la norma ISO 9000, enfocada en la gestión de la calidad; el
modelo por capas, que organiza los sistemas de forma modular; la CNCF, que fomenta la
adopción de tecnologías nativas de la nube; el Design Thinking, que sitúa al usuario en el centro
del proceso de innovación; TOGAF, que ayuda a alinear la estrategia empresarial con los
procesos de TI; y la norma ISO/IEC 25010, que evalúa la calidad del software y los sistemas.

\section{PMBOK}
En el ámbito de la gestión de proyectos y el desarrollo de software, es esencial contar con marcos
de referencia que ofrezcan una estructura robusta y metodologías efectivas para enfrentar los
retos actuales. Este proyecto se fundamenta en varios marcos de referencia reconocidos a nivel
internacional, cada uno con su enfoque particular y conjunto de principios que orientan las
prácticas organizacionales y técnicas. Entre estos se incluyen el PMBOK, que establece buenas
prácticas en la gestión de proyectos; la norma ISO 9000, enfocada en la gestión de la calidad; el
modelo por capas, que organiza los sistemas de forma modular; la CNCF, que fomenta la
adopción de tecnologías nativas de la nube; el Design Thinking, que sitúa al usuario en el centro
del proceso de innovación; TOGAF, que ayuda a alinear la estrategia empresarial con los
procesos de TI; y la norma ISO/IEC 25010, que evalúa la calidad del software y los sistemas.



\section{ISO 9000}
Según la Organización Internacional de Normalización, la serie de normas ISO 9000 se refiere a
un conjunto de principios y directrices internacionales que permiten a las organizaciones
satisfacer las expectativas y necesidades de sus clientes de manera consistente (Porfert, 1986).
Estas normas no solo ofrecen un marco sistemático para la mejora continua de los procesos, sino
que también garantizan la eficiencia operativa y la calidad de los productos y servicios ofrecidos.
Por lo tanto, pueden ser utilizadas para evaluar y mejorar cualquier aspecto de una organización,
incluyendo liderazgo, estrategia y planificación, personal, información y conocimiento,
seguridad, prestación de servicios, calidad del producto y resultados finales (Gray et al., 2022).


\section{Modelo por capas}
El modelo por capas, conocido también como arquitectura en capas, se refiere a una estructura
de diseño en la que el sistema se divide en diferentes niveles o capas, cada una con una función
específica e independiente de las demás. Spray (2023) señala que este enfoque permite gestionar
mejor la complejidad del sistema, favoreciendo la separación de responsabilidades y facilitando
el mantenimiento y la evolución del software a lo largo del tiempo. Esto hace que el modelo por
capas ofrezca ventajas en términos de modularidad, permitiendo que los componentes de una
capa sean modificados o actualizados sin afectar directamente a las otras. Además, esta
arquitectura no solo contribuye a la reutilización y mejora de la escalabilidad del sistema, sino
que también aumenta la capacidad de gestionar cambios y la interoperabilidad entre los
componentes.


\section{TOGAF}
TOGAF (The Open Group Architecture Framework) se ha convertido en uno de los marcos de
referencia más utilizados para el desarrollo y gestión de arquitecturas empresariales. Este marco
proporciona un enfoque estructurado que permite a las organizaciones alinear la estrategia de
negocio con los procesos de TI, facilitando la toma de decisiones y reduciendo el uso de recursos
(Mumtaza et al., 2025). Su arquitectura por fases, que abarca la planificación, diseño,
implementación y monitoreo, no solo facilita la integración de sistemas y tecnologías, sino que
también permite flexibilidad y adaptabilidad en un entorno empresarial en constante evolución.

\section{ISO/IEC 25010}
ISO/IEC 25010 es un estándar internacional que establece un modelo de calidad del software y
sistemas, ampliamente utilizado para evaluar la calidad de productos y servicios de TI. Según lo
explica ISO (ISO/IEC 25010, 2011), este modelo define características como funcionalidad ,
eficiencia, seguridad, mantenibilidad y usabilidad, proporcionando un marco integral para
asegurar que los sistemas cumplan con los requisitos del usuario y del negocio. La
implementación de ISO/IEC 25010 permite a las organizaciones no solo mejorar la calidad
interna y externa de sus productos, sino también garantizar que el software desarrollado sea
adaptable, seguro y eficaz en entornos dinámicos y complejos.

\section{HTCondor}
Como está expuesto en su página web (HTCondor, s.f), este es un sistema
especializado de gestión de cargas de trabajo para tareas de cómputo intensivo. Al igual que otros
sistemas de procesamiento por lotes completos, HTCondor proporciona un mecanismo de
encolado de trabajos, una política de planificación, un esquema de prioridades, monitoreo de
recursos y gestión de recursos. Los usuarios envían sus trabajos, ya sean seriales o paralelos, a
HTCondor; HTCondor los coloca en una cola, elige cuándo y dónde ejecutarlos según una
política, monitorea cuidadosamente su progreso y, finalmente, informa al usuario una vez que se
han completado.

\section{Universo HTCondor}
Así como está expresado en su página oficial (HTCondor, s.f), un universo HTCondor
está definido como aquel ambiente de ejecución en el cual una tarea es procesada (HTCondor,
n.d.). Los universos existentes a la fecha son los siguientes: vanilla, grid, java, scheduler, local,
parallel, vm, container y Docker.

\section{Universo GRID}
El universo GRID en HTCondor constituye un mecanismo especializado de ejecución de tareas 
computacionales que permite el envío de trabajos a recursos distribuidos ubicados en diferentes 
sistemas de cómputo, incluyendo servicios de computación en la nube, sistemas de gestión de lotes 
remotos y otras infraestructuras de grid computing. Este universo representa una abstracción fundamental
que extiende las capacidades de HTCondor más allá de los pools locales, habilitando la interoperabilidad
con múltiples plataformas heterogéneas de cómputo distribuido.